\section{Discussion}
\label{sec:discussion}

\todo{rewrite after the numbers land.}

\paragraph{Temporal context matters.}
Reactive baselines only see the current query or the current action,
so a slow-drifting attack (MINJA's gradient-descent flavor) can walk
just under their threshold. Temporal monitoring reframes detection as
``is this agent now different from itself six steps ago?'' -- a
question that is easier for the defender to answer than ``is this
query malicious?''.

\paragraph{Pivot A (\texorpdfstring{$\awt{} = 0$}{AWT = 0}).}
In the original proposal the detector was expected to produce
strictly positive Advance Warning Time. In practice the sharper
formulation is \emph{concurrent} detection at the moment the attack
takes effect; that is what the Phase 10 runs measure. This pivot is
documented in \texttt{docs/phase1\_decision.md}.

\paragraph{Limitations.}
\begin{itemize}
  \item The signal set is fixed at six scalars; an attacker with
  access to the signal catalog could in principle optimize a payload
  that masks every channel.
  \item Chronos (T5-small) is a zero-shot forecaster; domain adaptation
  to individual agent behavior is future work.
  \item The evaluation is limited to two published attack recipes.
  Unknown attacks are not covered; \sys{} is a complement, not a
  replacement, for input-level guards.
  \item Our throughput numbers come from the MockBackend simulator;
  real-LLM numbers land with the Production Hardening phase (Phase 9).
\end{itemize}

\paragraph{Threat model boundary.}
\sys{} defends against an attacker who can write to shared memory but
does not have direct control over an agent's weights or its system
prompt. Weight-level compromises and prompt-level compromises are out
of scope and would be flagged by orthogonal defenses.
